% titel: Blatt 2, A7: Erklärung der Grundbegriffe
% Teil II: die erarbeiteten Erklärungen.

\subsection*{Gegeben und gesucht}
Fünf chemische Begriffe fachsprachlich korrekt erklären.

\subsection*{Erklaerung}
\begin{itemize}
  \item[a)] \textbf{Edukt / Produkt:} Edukte sind die Ausgangsstoffe einer
        Reaktion (linke Gleichungsseite), Produkte die dabei gebildeten Stoffe
        (rechte Seite). Edukte werden im Lauf der Reaktion verbraucht,
        Produkte aufgebaut.
  \item[b)] \textbf{Stochiometrische Koeffizienten:} die ganzzahligen Zahlen
        vor den Formeln in einer ausbalancierten Reaktionsgleichung; sie
        geben an, in welchem Stoffmengenverhaeltnis die Stoffe reagieren.
  \item[c)] \textbf{$l$, $s$, $aq$, $g$:} Zustandssymbole --- $l$ = flüssig,
        $s$ = fest, $aq$ = in Wasser gelöst, $g$ = gasförmig.
  \item[d)] \textbf{Begrenzender Reaktand:} derjenige Edukt-Stoff, dessen
        Stoffmenge als erstes aufgebraucht ist; er begrenzt die maximale
        Produktmenge. Die anderen Reaktanden sind im Ueberschuss.
  \item[e)] \textbf{tatsaechliche vs.\ theoretische Ausbeute:} die
        theoretische Ausbeute ist die maximal mögliche Produktmenge laut
        Stochiometrie; die tatsaechliche ist die real erhaltene Menge. Lauft
        die Reaktion unvollstaendig oder treten Nebenreaktionen auf, ist die
        tatsaechliche Ausbeute kleiner; das Verhaeltnis in Prozent ist die
        Ausbeute (z.~B. $0{,}8\ \mathrm{mol}$ von $1\ \mathrm{mol}$
        theoretisch = $80\%$).
\end{itemize}

\begin{ergebnis}
\begin{itemize}
  \item[] Edukt = Ausgangsstoff, Produkt = gebildeter Stoff.
  \item[] Koeffizienten = Mengenverhaeltnisse der Formeln.
  \item[] $l$ flüssig, $s$ fest, $aq$ in Wasser gelöst, $g$ gasförmig.
  \item[] Begrenzender Reaktand = zuerst verbrauchter Edukt, bestimmt die
        Ausbeute.
  \item[] Theoretische Ausbeute = maximal möglich, tatsaechlich = real
        gewonnen ($\le$ theoretisch).
\end{itemize}
\end{ergebnis}

\begin{kontrolle}
Wortlaut gegen die Definitionen aus dem Skript / pro Sicherheit mit einem
konkreten Beispiel verifizieren (z.~B. $2\,\mathrm{H_2}+\mathrm{O_2}$:
ist $1\ \mathrm{mol}\ \mathrm{H_2}$ und $1\ \mathrm{mol}\ \mathrm{O_2}$
gegeben, begrenzt $\mathrm{H_2}$, denn $1\ \mathrm{mol}\ \mathrm{H_2}$
braucht nur $0{,}5\ \mathrm{mol}\ \mathrm{O_2}$).
\end{kontrolle}

\begin{merksatz}
Jeder Rechenbegriff haengt an der Stochiometrie: Koeffizienten liefern
Verhaeltnisse, der kleinste Reaktand bestimmt die Ausbeute, und die
Zustandssymbole gehören zur vollstaendigen Gleichung.
\end{merksatz}
